\section{Proofs}

\subsection{Proof of Proposition \ref{prop:axm-research-automation}}

The ratio of the conditional demands in \eqref{eq:axm-research-demands} is
\[
 \frac{AM}{H}
 =\left(\frac{\omega_M}{\omega_H}\right)^{\sigma_{HM}}
  (Aw)^{\sigma_{HM}}.
\]
Because \(s_M/(1-s_M)=M/(wH)=(AM/H)/(Aw)\), substitution yields
\[
 \frac{s_M}{1-s_M}
 =\left(\frac{\omega_M}{\omega_H}\right)^{\sigma_{HM}}
  (Aw)^{\sigma_{HM}-1}.
\]
Taking the stated limit proves all three cases. \(\square\)

\subsection{Proof of Proposition \ref{prop:axm-complements}}

Normalize the monopoly condition by \((K/L)^\alpha\) and write \(x=X/L\).
The normalized marginal cost tends to zero by hypothesis. At an interior
monopoly allocation,
\[
 p_X\left[1-\frac{1-s_X}{\sigma_{XL}}-\alpha s_X\right]=\frac{1}{A}.
\]
On the regular branch, the zero-marginal-cost limit sets the term in square
brackets to zero. Solving gives
\[
 \bar s_X=\frac{1-\sigma_{XL}}{1-\alpha\sigma_{XL}}.
\]
For \(\sigma_{XL}<1\), this number is strictly between zero and one. The CES share
mapping \eqref{eq:axm-sx} is one-to-one in \(X/L\), so \(X/L\) converges to a
finite positive constant. Hence \(Z/L\) converges to a constant and
\(Y\) is asymptotically a constant multiple of \(K^\alpha L^{1-\alpha}\).

Now suppose, as in the proposition, that a finite-rate balanced path exists.
The household Euler equation requires a constant \(Y/K\), while constant
production shares require \(K/L\) to be constant. Since \(L\) grows at \(n\),
\(K,Y\), and \(C\) grow at \(n\). Output per person and the wage are therefore
constant. \(\square\)

\subsection{Proof of Proposition \ref{prop:axm-cd-bgp}}

With \(\sigma_{XL}=1\), the final technology is
\[
 Y=K^\alpha L^\lambda X^\beta.
\]
The inverse demand elasticity is \(1-\beta\), so the monopoly first-order
condition becomes \(\beta p_X=1/A\). Since \(p_XX=\beta Y\),
\(X/A=\beta^2Y\), proving \eqref{eq:axm-inference-share-cd}. Replacing
\(X\) by \(\beta^2AY\) gives \eqref{eq:axm-reduced-cd}. Taking growth rates
and using \(g_K=g_Y\) and \(g_L=n\) yields
\[
 g_Y=n+\nu g_A.
\]

On a balanced path, raw research compute has a constant output share, so
\(g_M=g_Y=n+\nu g_A\). If \(\sigma_{HM}=1\), human research grows at \(n\) and
\[
 g_E=\omega_H n+\omega_M(g_A+g_M)
 =n+\omega_M(1+\nu)g_A.
\]
If \(\sigma_{HM}>1\) and \(s_M\to1\), then \(E\) is asymptotically proportional
to \(AM\), giving
\[
 g_E=g_A+g_M=n+(1+\nu)g_A.
\]
Both cases are \(g_E=n+\bar s(1+\nu)g_A\). Since
\(\dot A=\chi E^\eta\) and a constant \(g_A>0\) requires
\(g_{\dot A}=g_A\),
\[
 g_A=\eta[n+\bar s(1+\nu)g_A].
\]
Solving gives \eqref{eq:axm-cd-growth-rates}; a positive solution requires
\(D(\bar s)>0\).

The Euler equation gives \(K/Y=\alpha/(\rho+\delta+\gamma)\), which implies the
investment share in \eqref{eq:axm-cd-shares}. Let \(m=M/Y\) and
\(u=\beta^2\). The research-compute first-order condition can be written
\[
 q\,\eta\bar s\,\frac{\dot A}{M}=1,
 \qquad\text{so}\qquad
 \frac{qA}{Y}=\frac{m}{\eta\bar s g_A}.
\]
Because this ratio is constant, \(g_q=g_Y-g_A\). The envelope derivative of
operating profits implies
\[
 \frac{\pi_A}{q}
 =\frac{uY/A}{q}
 =\frac{u\eta\bar s g_A}{m}.
\]
Substitute these expressions into the costate equation and use
\(g_Y=n+\gamma\) and the Euler condition:
\[
 n+\gamma-g_A
 =\rho+\gamma-\eta\bar s g_A
  -\frac{u\eta\bar s g_A}{m}.
\]
Solving for \(m\) gives \eqref{eq:axm-cd-shares}. The resource constraint gives
\(c=1-u-i-m\). Positivity, interiority, second-order, and transversality
conditions are separate requirements, as stated. \(\square\)

\subsection{Proof of Proposition \ref{prop:axm-no-bgp}}

Suppose for contradiction that an equilibrium satisfying the hypotheses has
finite constant growth rates. With \(\sigma_{XL}>1\), AI services dominate the
CES as their relative quantity becomes large. The limiting final technology is
proportional to \(K^\alpha X^{1-\alpha}\). The inverse demand elasticity tends
to \(\alpha\), and the monopoly first-order condition implies
\[
 \frac{X}{A}=(1-\alpha)^2Y[1+o(1)].
\]
Substituting \(X\) back into final production gives
\[
 \frac{Y}{K}=\mathcal D A^{(1-\alpha)/\alpha}[1+o(1)]
\]
for \(\mathcal D>0\). Since \(A\to\infty\), \(Y/K\to\infty\). The gross return
to capital \(\alpha Y/K\) therefore diverges. Equation \eqref{eq:axm-euler}
then makes consumption growth unbounded, contradicting finite constant growth.
\(\square\)

\subsection{Derivation of the AI-dominated singular limit}

Let \(z=Y/K\), \(u=U/Y\), \(m=M/Y\), \(i=(\dot K+\delta K)/Y\), and
suppose the limiting shares are positive and constant. Final-production and
monopoly optimality give
\[
 u\to(1-\alpha)^2,\qquad z\sim\mathcal D A^\theta,\qquad
 \theta=\frac{1-\alpha}{\alpha}.
\]
Write \(g_A\sim hz\). The Euler equation and a constant consumption share imply
\[
 \frac{g_Y}{z}\to\alpha.
\]
Since \(g_Y=g_K+g_z\), \(g_K/z\to i\), and \(g_z=\theta g_A\),
\begin{equation}
 i+\theta h=\alpha.
 \label{eq:axm-app-resource-growth}
\end{equation}

Automated research implies \(\dot A\sim\chi(AM)^\eta\). Log-differentiating
\(g_A=\dot A/A\), using \(g_M=g_Y\) for constant \(m\), gives
\[
 g_{g_A}=(\eta-1)g_A+\eta g_Y.
\]
But \(g_A=hz\) implies \(g_{g_A}=g_z=\theta g_A\). Divide by \(z\):
\[
 \theta h=(\eta-1)h+\eta\alpha,
\]
so
\[
 h=\frac{\eta\alpha}{1+\theta-\eta}.
\]
Equation \eqref{eq:axm-app-resource-growth} gives \(i=\alpha-\theta h\).

The research-compute first-order condition implies
\[
 \frac{qA}{K}=\frac{m}{\eta h}.
\]
Operating-profit differentiation gives
\(\pi_A/q=uz/(qA/K)\). A constant \(qA/K\) means
\(g_q=g_K-g_A\). Substituting these results into the costate equation and using
\eqref{eq:axm-app-resource-growth} yields
\[
 \frac{qA}{K}=\frac{u}{\eta\alpha},
 \qquad
 m=\frac{\eta u}{1+\theta-\eta}.
\]
The resource constraint determines \(c\). For \(0<\eta<1\), both \(i\) and \(c\)
are positive; in fact
\[
 c=\alpha(1-\alpha)
 \left(1+\frac{\eta}{1+\theta-\eta}\right)>0.
\]
Finally, \(g_z=\theta g_A\sim\theta hz\), so
\(\dot z\sim\theta hz^2\), which has a finite blow-up time.

For the wage, the CES share equation gives
\[
 1-s_X\asymp(L/X)^{(\sigma_{XL}-1)/\sigma_{XL}}.
\]
Using \(w=(1-\alpha)(1-s_X)Y/L\), \(X\asymp AY\), and ignoring the bounded
population growth term relative to \(z\), we obtain
\[
 \frac{g_w}{z}\to
 \frac{\alpha-(\sigma_{XL}-1)h}{\sigma_{XL}}.
\]
This expression is positive when
\(\sigma_{XL}<1/(\alpha\eta)\). Because \(\eta<1\),
the wage limit has the sign stated in the conditional result. \(\square\)

\section{Numerical equilibrium audit}

For every reported path, the code recomputes the static allocation at each date
and records the following residuals: the goods-market identity, the labor-market
identity, the monopoly first-order condition, both research first-order
conditions, the household Euler equation, and the four differential equations.
The collocation derivative, rather than a plotted finite difference, is used to
audit the dynamic equations. The simulation files and their residual summaries
are stored in \texttt{numerical\_axm}; the independent solver is
\texttt{scripts/simulate\_axm\_equilibrium.py}, with a separate free-boundary
solver for the gross-substitutes regime.
